% Source: Weinberg GR, physical PDF pp. 87--89
%         (printed pp. 61--63).
% Coverage: Section 2.13, Temporal Order and Antiparticles; equations
%           (2.13.1)--(2.13.2).
% Figures/tables/footnotes: optional-reading footnote; reference marker 13.
% Status: source-reviewed and compile-clean.
% Uncertainties: none.

\subsection{Temporal Order and Antiparticles\optionalreading}
\label{sec:2.13}
\begingroup
\renewcommand\thefootnote{*}
\footnotetext{This section lies somewhat out of the book's main line of
development, and may be omitted in a first reading.}
\endgroup

One of the most striking features of the Lorentz transformations is that they
do not leave invariant the order of events. For instance, suppose that in one
reference frame an event at \(x_2\) is observed to occur later than one at
\(x_1\); that is, \(x_2^0>x_1^0\). A second observer who sees the first
observer moving with velocity \(\mathbf v\) will see the events separated by a
time difference
\[
  x_2'^0-x_1'^0
  =\Lambda^0{}_\mu(\mathbf v)(x_2^\mu-x_1^\mu),
\]
where \(\Lambda^\mu{}_\nu(\mathbf v)\) is the ``boost'' defined by
Eqs.~\eqref{eq:2.1.17}--\eqref{eq:2.1.21}. Applying
Eqs.~\eqref{eq:2.1.17} and \eqref{eq:2.1.21} gives then
\[
  x_2'^0-x_1'^0
  =\gamma(x_2^0-x_1^0)
   +\gamma\mathbf v\cdot(\mathbf x_2-\mathbf x_1),
\]
and this will be negative if
\begin{equation}
  \mathbf v\cdot(\mathbf x_2-\mathbf x_1)
  <-(x_2^0-x_1^0).
  \tag{2.13.1}\label{eq:2.13.1}
\end{equation}
At first sight this might seem to raise the danger of a logical paradox.
Suppose that the first observer sees a radioactive decay
\(A\to B+C\) at \(x_1\), followed at \(x_2\) by absorption of particle
\(B\), for example, \(B+D\to E\). Does the second observer then see \(B\)
absorbed at \(x_2\) before it is emitted at \(x_1\)? The paradox disappears if
we note that the speed \(\lvert\mathbf v\rvert\) characterizing any Lorentz
transformation \(\Lambda(\mathbf v)\) must be less than unity, so that
Eq.~\eqref{eq:2.13.1} can be satisfied only if
\begin{equation}
  \lvert\mathbf x_2-\mathbf x_1\rvert
  >\lvert x_2^0-x_1^0\rvert .
  \tag{2.13.2}\label{eq:2.13.2}
\end{equation}
However, this is impossible, because particle \(B\) was assumed to travel from
\(x_1\) to \(x_2\), and Eq.~\eqref{eq:2.13.2} would require its speed to be
greater than unity; that is, than the speed of light. To put it another way,
the temporal order of events at \(x_1\) and \(x_2\) is affected by Lorentz
transformations only if \(x_1-x_2\) is \emph{spacelike}; that is,
\[
  \eta_{\mu\nu}(x_1-x_2)^\mu(x_1-x_2)^\nu>0,
\]
whereas a particle can travel from \(x_1\) to \(x_2\) only if \(x_1-x_2\) is
\emph{timelike}; that is,
\[
  \eta_{\mu\nu}(x_1-x_2)^\mu(x_1-x_2)^\nu<0 .
\]

Although the relativity of temporal order raises no problems for classical
physics, it plays a profound role in quantum theories. The uncertainty
principle tells us that when we specify that a particle is at position
\(\mathbf x_1\) at time \(t_1\), we cannot also define its velocity precisely.
In consequence there is a certain chance of a particle getting from \(x_1\) to
\(x_2\) even if \(x_1-x_2\) is spacelike; that is,
\(\lvert\mathbf x_1-\mathbf x_2\rvert
>\lvert x_1^0-x_2^0\rvert\). To be more precise, the probability of a
particle reaching \(x_2\) if it starts at \(x_1\) is nonnegligible as long as
\[
  \lvert\mathbf x_1-\mathbf x_2\rvert^2
  -(x_1^0-x_2^0)^2\lesssim\frac{\hbar^2}{m^2},
\]
where \(\hbar\) is Planck's constant (divided by \(2\pi\)) and \(m\) is the
particle mass. (Such space-time intervals are very small even for elementary
particle masses; for instance, if \(m\) is the mass of a proton then
\(\hbar/m=2\times10^{-14}\) cm, or in time units
\(6\times10^{-25}\) sec. Recall that in our units
\(1\) sec \(=3\times10^{10}\) cm.) We are thus faced again with our paradox;
if one observer sees a particle emitted at \(x_1\), and absorbed at \(x_2\),
and if
\(\lvert\mathbf x_1-\mathbf x_2\rvert^2-(x_1^0-x_2^0)^2\) is positive
(but less than \(\hbar^2/m^2\)), then a second observer may see the particle
absorbed at \(x_2\) at a time \(t_2\) before the time \(t_1\) it is emitted at
\(x_1\).

There is only one known way out of this paradox. The second observer must see a
particle emitted at \(x_2\), and absorbed at \(x_1\). But in general the
particle seen by the second observer will then necessarily be different from
that seen by the first. For instance, if the first observer sees a proton turn
into a neutron and a \emph{positive} pi-meson at \(x_1\), and then sees the
pi-meson and some other neutron turn into a proton at \(x_2\), then the second
observer must see the neutron at \(x_2\) turn into a proton and a particle of
\emph{negative} charge, which is then absorbed by a proton at \(x_1\) that
turns into a neutron. Since mass is a Lorentz invariant, the mass of the
negative particle seen by the second observer will be equal to that of the
positive pi-meson seen by the first observer. There is such a particle, called
a negative pi-meson, and it does indeed have the same mass as the positive
pi-meson. This reasoning leads us to the conclusion that for every type of
charged particle there is an oppositely charged particle of equal mass, called
its antiparticle. Note that this conclusion does not obtain in nonrelativistic
quantum mechanics or in relativistic classical mechanics; it is only in
relativistic quantum mechanics that antiparticles are a
necessity.\textsuperscript{\hyperref[ch2-ref-13]{13}} And it is the existence
of antiparticles that leads to the characteristic feature of relativistic
quantum dynamics, that given enough energy we can create arbitrary numbers of
particles and their antiparticles.
